\section{Constuction of Boneh-Franklin IBE with Chosen Ciphertext Security}

In his original paper \cite{Shamir41} Shamir introduced a public key encryption
scheme that resembled an ideal mail system, in which the public key can be an
arbitrary string which  uniquely identifies the user in a way that he/she
cannot later deny and that is readily available to the other party.\cite{Shamir41}

Shamir's original motivation for identity-based encryption was to simplify the
certificate, signing, and key management processes with regards
for example in a mail system.\cite{BF01}.

Let's say when \textit{Alice} wants to send  a mail to \textit{Bob} at \url{bob@company.com}
she simply encrypts her message using the public key string \lstinline{bob@company.com}. 
In the proposed cryptosystem there is no need for \textit{Alice} to obtain \textit{Bob}’s public key
certificate. When Bob getting noticed that he received the encrypted mail he contacts a trusted
third party called the \textit{Private Key Generator} (denoted as \textit{PKG}  furthermore) and 
after authenticating himself to the \textit{PKG} \textit{Bob} obtains his private key
from the \textit{PKG}. \cite{BF01}

The authentication takes places in the same way \textit{Bob} would authenticate
himself to a trusted Certificate Authority (denoted as CA furthermore).
\textit{Bob} can then therefore read his e-mail sent by \textit{Alice}.

In comparison to an existing secure e-mail infrastructure,
with the use of the proposed fully functional
identity-based encryption scheme\cite{BF01}
\textit{Alice} can send an encrypted mail to \textit{Bob} even if he has not yet
setup his public key certificate. Regarding the system's performance,
now it becames comparable to the performance of
ElGamal encryption\cite{cryptosys.elgamal} in $\mathcal{F}^{*}_p$. The security of the
Boneh-Frenkel scheme is analogues to the computational Diffie-Hellman assumption\cite{crypto.cdhddh},
namely the new system has chosen ciphertext security in the random oracle model\cite{Bellare1993}.

Also the authors showed that in their scheme \textit{PKG}
can be distributed using standard techniques from treshold
cryptography\cite{crypto.seckeydistdl}, consequently given the
\textit{Boneh-Frenkel} scheme based IBE, the
\lstinline{master-key} is never available in a single location.

%% TODO:    itt az implemnetacio-t be kell hivatkozni +
%%          az architekturat ennek megfeleloen kell oszerakni
%%          docker-ben. --> lasd kesobb

% Fujisaki-Okamoto [16]

%% TODO:    Detailed building blokk --> /appendix/{}
\subsection{Building blocks for Boneh-Franklin IBE Scheme}
 \label{ssec:ibe.bonehfranklin.scheme}

% Choosing the right Curve families

%% [] <- ehhez meg kell az ETSI is
%%  TODO: BF01-ben nincs benne, hogy melyik curve-re van illesztve. --> ez appendixbe
%%  TODO: A valasztott gorbe alapjan kell erre majd a unit teszteket lefuttattni.
%%      TODO:   Ha van ido, akkor tobb cure parameterrre is ki kell terjeszteni a parameter tesztet.

% Curve
The following table summarizes some curve families and their repsective parameters
based on \cite{ecc.fam.param}  with regard to \textit{pairing-friendly elliptic curves}\cite{ecc.pairings.param}
(for a detailed summarization regarding elliptic curve theory,
please refer to \textit{Appendix Sections
\ref{sec:apx.ecc.curves}} and \textit{\ref{sec:apx.ecc.curves.pairing}})

%
\begin{table}[h!]
\centering
 \begin{tabular}{||c c c c||}
 \hline
 Col1 & Col2 & Col2 & Col3 \\ [0.5ex]
 \hline\hline
 1 & 6 & 87837 & 787 \\
 2 & 7 & 78 & 5415 \\
 3 & 545 & 778 & 7507 \\
 4 & 545 & 18744 & 7560 \\
 5 & 88 & 788 & 6344 \\ [1ex]
 \hline
 \end{tabular}
\end{table}

% =======================================
% Curve paremeters (\cite{ecc.fam.param})
%% TODO: safe_curve_parameters from \cite{ecc.sec.param}

% k := embedding degree
% \rho = balance (\rho = 1 is the optimal)
% (prime) p(x) :=
% (prime) r(x) :=
% t(x) :=
% r := curve order
% ---

% BN (Barreto-Naehrig) curves (used to be optimal at the 128-bit security level):
% k = 12,
% \rho \approx 1
% p(x) = 36x^4 + 36x^3 + 24x^2 + 6x + 1
% r(x) = 36x^4 + 36x^3 + 24x^2 + 6x + 1
% t(x) = 6z^2 + 1

% BLS12 curves:
% k = 12
% \rho \approx 1.5
% p(x) = (x-1)^2 (x^4 - x^2 +1)/3+x
% r(x) = x^4 - x^2 +1
% t(x) = x + 1

% BLS24 curves:
% k = 24
% \rho \approx 1.25
% p(x) = (x - 1)^2 (x^8 - x^4 +1)/3+x
% r(x) = x^8 -x^4 +1
% t(x) = x+1

% KSS18 curves:
% k = 18
% \rho \approx 4/3
% p(x) = (x^8 + 5x^7 + 7x^6+37x^5+188x^4 + 259 x^3 + 343 x^2 + 1763 x + 2401)/21
% r(x) = (x^6 + 37 x^3 + 343)/343
% t(x) = (x ^4 + 16x +7)/7

% ...







%% Weil pairing
Let  $\mathfrak{P} = (P) - (\infty)$ and $\mathfrak{Q} = (Q) - (\infty)$
be the divisors respectively, then for \textit{ Weil pairing}
the map is follows\cite{ecc.fam.param}:

%%% Consequently, functions can be evaulatiod over points istead of divisors:
\[
    \begin{matrix}
        w & : & E(\mathcal{F}_q)[r]  \times E(\mathcal{F}_{q^k})[r] \rightarrow \mathcal{F}^{*}_{q^k}  \\
        w_r(P, Q) & = & (-1)^r \frac{f_{r,P}(\mathfrak{Q}))}{f_{r,Q}(\mathfrak{P})}
    \end{matrix}

\]
%%%% -> RELIC: UNICAMP, flexible and state-of-the-art.



% =======================================







% Bilinear mapping Weil pairing
\subsubsection{Bilinear mapping and Weil pairing}

%% Pairing-based cryptgraphy
%%% asecuritysite: basics of pairing-cryptography <- using MIRACL

Let $G_1, G_2$ be abelian groups, written additively.
Now let $G_1$ be a multiplicative cyclic group of oder $n$.

Then, we can write the pairing as follows: $E: G_1 \times G_2 \rightarrow G_t$.

%%%%% Bilinearity cases <- ekvivalens atalakitasok.
%%%%%   TODO:   Erre kell 1 kulon unit teszt, mutatva, hogy az G_1, G_2-nel fennallnak a relaciok
%%%%%   TODO:   Itt a jeloles technikat ossze kell hozni a jegyzettel + a videokkal,
%%%%%           A vegleges jeloles legyen ugyanaz, mint a jegyzetben !

\begin{itemize}
    \item $e(P, S+T) = e(P,S) e(P(T))$, where $P \in G_1, S,T \in G_2$
    \item $e(aP, bQ) = e(bP, aQ)$,  where $a,b \in Z_n$ and $P \in G_1$, $Q \in G_2$
    \item $e(aP, bQ) = e(abP,Q)$,  where $a,b \in Z_n$ and $P \in G_1$, $Q \in G_2$
    \item $e(aP, bQ) = e(P,Q)^{ab}$,  where $a,b \in Z_n$ and $P \in G_1$, $Q \in G_2$
    \item $e(-P, Q) = e(P,Q)^{-1} = e(P,-Q)$, where $P \in G_1, Q \in G_2$
    \item $e(P,0) = e(0, Q) = 1$ , where $P \in G_1, Q \in G_2$
    \end{itemize}

% TODO: Alap implementaciohoz: 1 hoston minden.
% TODO: A jeloleseket
\subsection{IBE Boneh Franklin scheme (simplified)}

Overview IBE Boneh-Franklin scheme based on \cite{ecc.fam.param}
%% Jelolesek
\begin{itemize}
    \item \textbf{Initialization phase:}
    \item   \begin{enumerate}
                \item \textit{PKG} generates $\textit{master-key}$ $s\in \mathcal{Z}^{*}_{r}$ and computes public key $P_{pub} = sP$
                \item Fix hash functions as follows: $H_1: \{ 0,1 \}^* \rightarrow \{ \mathcal{G}_1\}$ and $H_2: \{ \mathcal{G}_T\} \rightarrow \{0,1\}^m$
            \end{enumerate}
    \item \textbf{Key generation phase:}
    \item   \begin{enumerate}
                \item User with identity $ID_i$ computes $P_i$ public key as: $P_i = H_1(ID_i)$
                \item \textit{PKG} generates private key $S_i = sP_i$
    \end{enumerate}
    \item \textbf{Encryption phase:}
    \item   \begin{enumerate}
                \item To encrypt $m$, \textit{Bob} selects randomly $l$ and computes $R=lP$ and $c$ ciphertext:
                \begin{equation*} c = m \oplus H_2(e(P_A, P_{pub})^l ) \end{equation*}
                \item \textit{Bob} sends $(R,c)$ to \textit{Alice}.
    \end{enumerate}
    \item \textbf{Decryption phase:}
    \item   \begin{enumerate}
                \item \textit{Alice} uses her own $S_A$ private key to compute $m$ as:
                \begin{eqnarray*} c \oplus H_2(e(S_A, R)) = c \oplus H_2(e(sP_A, lP)) = c \oplus H_2(e(P_A, P_{pub})^l)=m  \end{eqnarray*}
    \end{enumerate}
\end{itemize}


\subsection{Attacker Model}

In comparison to standard models for chosen ciphertext security
    %% cite: 37, 2

within the Boneh-Franklin scheme the \textit{Adversary}
is more powerful:

    %%  Az alap tamadasi modell-ben tetszologes meritest vegez az adversary
    %%  a private kulcsok-ban, majd public key ID-ket tamad
    %%  Magyaran:   Obtaineli a PKG-tol H(master_key)-t majd a Hash-en
    %%              collusion-t keres public ID-k alapjan
    %% A security-t a Dlog assumption adja:
    %%  G1 gorbe-n leve S(ecret) pont es
    %%  G2 gorbe-n levo P(ublic) pont-kat feleltessuk meg
    %%  G3(U,V)-ra projektalt gorbenkbol.

    Firstly, the attacker is allowed to attack any arbitrary \textit{public key ID}
    Secondly, while mounting  the choosen ciphertext attack on the \textit{ID}
    the attacker is allowed to obtain any private keys corresponding to some
    identites other than the private key for \textit{ID}. The last
    condition is equivalent to the case where the attacker obtains a number
    of private keys corresponding to some identites, than it tries to attack
    some other \textit{public key ID}.

    %%  Ez lemodellezheto egyszeruen halozat nelkul is:
    %%      ./volume-ra legyen kiirva 1 file-ba 1 rakas privatekulcs
    %%      Eve ferjen hozza a ./volume-n tarolt file-hoz es futtasson tesztet
    %%      ezek alapjan a kulcsok alapjan.
    %%      Ha lesz collision H(pub_key) H(lopott_priv_key) kozott akkor nyer.






% --------------------------------------